\section{Finite Haar constraints and connected lattice integrals}
\label{sec:middle-static}

Rounds 13--22 developed several distinct mathematical models. Throughout this
section, \(x_p=\Tr(U_p)/2\) denotes a normalized fundamental Wilson trace,
\(\chi_p=2x_p\) its ordinary character, and \(C_e\) the positive link Casimir
with eigenvalues \(j(j+1)\). A finite physical Hamiltonian has the convention
\[
 H=\alpha\sum_e C_e-\sum_p\lambda_p x_p,\qquad \alpha>0.
\]
Here \(\alpha,\lambda_p\) have energy units. Euclidean coefficients
\(\kappa_p\), introduced in densities proportional to
\(\exp(\sum_p\kappa_p x_p)\), are dimensionless and do not define a physical
clock. Every spectral statement below retains a fixed positive energy
reference \(E_\star\). The symbol \workstar\ identifies a result derived or
computed in the workbench; it does not assert priority over the literature.
The audit ledger accounts for all 62 historical entries in this range,
including four named Round13 studies and 58 numbered loops.

\subsection{Compact moments, response, and complete interval coverage \workstar}
\label{sec:middle-moments}

The one-holonomy model uses normalized Haar pushforward
\[
 d\mu_\kappa(x)=Z_\kappa^{-1}e^{\kappa x}
       \frac2\pi\sqrt{1-x^2}\,dx,\qquad -1\le x\le1.
\]
Writing \(m_n=\int x^n\,d\mu_\kappa\), integration of the derivative of
\(x^ne^{\kappa x}(1-x^2)^{3/2}\) gives the familiar identity
\begin{equation}
 nm_{n-1}-(n+3)m_{n+1}+\kappa(m_n-m_{n+2})=0.
 \label{eq:middle-moments}
\end{equation}
The first term is absent at \(n=0\). For nonzero \(\kappa\), normalization
and \(m_1=u\) express every moment affinely in \(u\). At zero coupling one
instead uses \(m_{2j}=C_j/4^j\), \(m_{2j+1}=0\), where \(C_j\) is a Catalan
number. Positivity supplies
\[
 (m_{i+j})_{i,j=0}^{r}\succeq0,\qquad
 (m_{i+j}-m_{i+j+2})_{i,j=0}^{r-1}\succeq0.
\]
The workbench implementation retains rational vectors whose quadratic forms
are affine inequalities in \(u\), so rejected means have exact polynomial
witnesses. Its 28 finite certificates include a level-six interval at
\(\kappa=1\) of width approximately \(1.40631\,10^{-13}\); the analogous
width at \(\kappa=5\) is approximately \(6.03663\,10^{-8}\).
These are enclosure widths, not sampling errors.

Completeness is a separate argument. All moment and support inequalities
produce a probability measure on \([-1,1]\), using multiplication by \(x\)
as a self-adjoint contraction on the polynomial completion. All recurrence
identities extend to \(C^1\) tests and imply
\[
 D[(1-x^2)\mu]=[\kappa(1-x^2)-3x]\mu.
\]
The interior density is proportional to
\(e^{\kappa x}\sqrt{1-x^2}\); the equation excludes endpoint atoms.
Consequently the full hierarchy identifies the desired measure. Nested
finite feasible intervals shrink to a point, without a proved rate.
Successful finite computations additionally require vanishing optimization
slack; the delivered level-six cap does not execute this limit.

Differentiation with respect to the coupling gives \(u'=\operatorname{Var}(x)>0\).
Combining this with \eqref{eq:middle-moments} yields
\[
 \kappa u'+3u=\kappa(1-u^2),\qquad u(0)=0,\quad u'(0)=\tfrac14.
\]
This is a static response equation. The regular solution is
\(u=I_2(\kappa)/I_1(\kappa)\); singular modified-Bessel branches violate
boundedness at the origin. Replacing \(3u\) by \(2u\) changes the prior to
uniform measure on \([-1,1]\), rather than improving the Haar approximation.
The project also derived a useful residual certificate: fix
\(0<\kappa_0\le\kappa\) and a differentiable approximation \(p\ge0\)
throughout \([\kappa_0,\kappa]\), with
\(r=p'-(1-p^2-3p/\kappa)\),
\begin{equation}
 |p(\kappa)-u(\kappa)|
 \le \left(\frac{\kappa_0}{\kappa}\right)^3
       |p(\kappa_0)-u(\kappa_0)|
 +\kappa^{-3}\int_{\kappa_0}^{\kappa}s^3|r(s)|\,ds.
 \qquad\text{\workstar}
 \label{eq:middle-response}
\end{equation}
It follows by an integrating factor, since the error damping contains
\(3/\kappa+u+p\ge3/\kappa\). A certified launch error therefore does not
bound a later floating trajectory unless the entire residual is enclosed.
The point-mass proposal \(m_n=(1/3)^n\) at \(\kappa=9/8\) passes positivity
and the first identity but has next-identity residual \(8/9\).
This supplies an exact counterexample to a zero-variance closure.

Rounds14--15 moved to the genuine two-holonomy density
\(Z^{-1}e^{x+y+\eta z}dU\,dV\), with
\(z=\Tr(UV)/2\). For \(C(\eta)=\operatorname{Cov}_\eta(x,y)\),
differentiating the normalized quotient gives
\[
 C'(\eta)=\mathbb E_\eta[
 (x-\mathbb E_\eta x)(y-\mathbb E_\eta y)(z-\mathbb E_\eta z)],
 \qquad |C'|\le2.
\]
The last inequality uses \(|z-\mathbb Ez|\le2\), Cauchy--Schwarz, and
\(\operatorname{Var}(x),\operatorname{Var}(y)\le1\).
The first loop's small-\(\eta\) estimate did not resolve the chosen point
\((1,1,1/4)\). Exact total-degree Taylor integration did.
For any bounded insertion and action \(|S|\le M<N+2\), the standard tail
\[
 R_N(M)=\frac{M^{N+1}}{(N+1)!}\frac1{1-M/(N+2)}
\]
bounds numerator and partition separately. Jensen gives \(Z\ge1\) when
the unweighted action mean vanishes; signed interval division retains all
endpoint combinations. Point intervals \([L,U]\) extend across a cell of
radius \(\rho\) as \([L-2\rho,U+2\rho]\). All eight initial midpoint
intervals were positive, while all eight cell bounds failed. Refining only
the failing cells produced 8, 16, then 21 cells, with a final minimum lower
bound approximately \(0.0003346525\) over the whole closed interval
\(\eta\in[1/8,1/4]\). Increasing already adequate point precision would not
have removed the geometric transport loss.

These results support certified static response calculations and provide
countermodels to incomplete closures. Their next application requires
shared-link identities for larger graphs, not an identification of \(\eta\)
with time or energy. Sources:
\repo{research/round13/moments/README.md},
\repo{research/round13/response/response.md},
\repo{research/round14/forward/loop2.md},
\repo{research/round15/forward/A2/result.md}.

\subsection{Closed-surface gluing and the first shared face \workstar}
\label{sec:middle-gluing}

On the closed cube there are eight vertices, twelve links and six faces.
For \(\chi_n\), the spin-\(n/2\) character of dimension \(d_n=n+1\), write
\(f_p=\sum_n a_{p,n}\chi_n\). Schur orthogonality forces equal labels on
adjacent faces. Each edge contributes \(d_n^{-1}\), each surviving vertex
index \(d_n\), giving the graph-specific application
\begin{equation}
 \int\prod_{p=1}^{6}f_p(U_p)\,dU
       =\sum_{n\ge0}d_n^{-4}\prod_{p=1}^{6}a_{p,n},
 \qquad
 \mathbb E_0\prod_{p=1}^{6}x_p=2^{-10}.
 \qquad\text{\workstar}
 \label{eq:middle-cube}
\end{equation}
The identity is immediate for finite character polynomials; analytic
weights require convergence or the retained Taylor remainder. If any one
face weight equals one, only label zero survives and the other five class
observables factorize. Thus five-face factorization coexists with a
nonzero six-face moment. It would be incorrect to use this as evidence for
independent face holonomies.

For \(e^{\kappa x}\),
\[
 a_n(\kappa)=\sum_{j\ge0}
 \frac{(n+1)\kappa^{n+2j}}
 {2^{n+2j}j!(n+1+j)!}.
\]
An insertion of one \(x_p\) differentiates that face's coefficient before
setting all couplings equal. It does not differentiate the common-coupling
partition six times. At \(\kappa=1/4\), total degree18 first met the
specified expectation precision; degree24 gave
\(\mathbb E\prod_px_p\simeq0.001023116359923\) and a partition excess over
the independent-face product exceeding \(2.4598180\,10^{-7}\).
The exact fractions, including the failed coarse degrees, remain in the
source certificate.

Two adjacent cubes have twelve vertices, twenty links and eleven faces.
Conditioning on their shared holonomy \(H\), each five-face disk contributes
\(\sum_n d_n^{-4}(\prod_{\mathrm{disk}}a_{p,n})\chi_n(H)\).
The actual integral is therefore
\begin{equation}
 \sum_{n,m,k}
 \frac{\prod_{p\in L}a_{p,n}}{d_n^4}
 \frac{\prod_{p\in R}a_{p,m}}{d_m^4}
 a_{s,k}\,N_{nm}^{k},
 \qquad
 N_{nm}^{k}=\mathbf1_{\{|n-m|\le k\le n+m,\ n+m+k\ {\rm even}\}}.
 \quad\text{\workstar}
 \label{eq:middle-two-cube}
\end{equation}
The shared-edge incidence is three, so the old one-surface single-label
rule is insufficient. Center parity makes the all-eleven distinct-trace
Haar moment zero. With the shared trace squared, both fusion channels
\(k=0,2\) contribute and
\(\mathbb E_0[(\prod_{\mathrm{outer}\ 10}x_p)x_s^2]=2^{-19}\).
Removing \(k=2\) halves the answer.

Keeping the same all-eleven-trace observable, Round16 compared outer
couplings \(1/8\) and shared coupling \(1/8\) against shared coupling zero.
The degree24 difference is enclosed near \(2.4303274916\,10^{-7}\),
with width about \(4.173\,10^{-22}\). The zero shared-action baseline is
nonzero: its leading numerator at common outer coupling \(t\) is
\(41t^5/(81\,2^{18})\). Deleting a coefficient is not deleting an insertion.
These finite integrals expose graph dependence and supply later trial
matrix elements; they provide no transfer Hamiltonian by themselves.
Sources: \repo{research/round15/forward/B1/derivation.md},
\repo{research/round15/forward/B2/derivation.md},
\repo{research/round16/forward/loop1/derivation.md},
\repo{research/round16/forward/loop2/derivation.md}.

\subsection{Recoupling, sufficient coordinates, and surrounding weights \workstar}
\label{sec:middle-conditional}

The four-cube \(2\times2\times1\) complex has 18 vertices, 33 links and
20 faces. Integrating its central link in the adjoint representation
requires the complete invariant projector on \((\mathbb R^3)^{\otimes4}\).
The three invariant tensors
\(B_1=\delta_{ab}\delta_{cd}\),
\(B_2=\delta_{ac}\delta_{bd}\),
\(B_3=\delta_{ad}\delta_{bc}\)
have Gram matrix with diagonal nine and off-diagonal three. The
Clebsch--Gordan decomposition
\(1\otimes1=0\oplus1\oplus2\) proves completeness, and Haar averaging gives
\(P=B(B^TB)^{-1}B^T\). The inverse has diagonal \(2/15\) and
off-diagonal \(-1/30\). For an explicitly realized tetrahedral boundary,
the product of four normalized adjoint traces has expectation
\(-1/405\); diagonal-only and one-channel substitutes instead give
\(2/405\) and \(1/729\). The negative entries cannot be dropped.

Writing the central quaternion as \(q\in S^3\), face traces as
\(q\cdot a_i\), and the action as \(q\cdot b\), with
\(b=\sum_i\kappa_i a_i\), the partition depends only on \(\norm b\).
The joint observable also depends on the directions \(a_i\). Tetrahedral
and commuting boundaries can have the same \(b\) but opposite-sign
joint expectations. The complete Gram matrix of
\((b,a_1,\ldots,a_4)\) is sufficient: equality of Grams defines an isometry
between the spans, extends to an element of \(O(4)\), and preserves Haar
sphere measure. This proof includes singular Grams and \(b=0\), without
inverting a Gram matrix. Admissibility requires positive semidefiniteness,
rank at most four, unit \(a_i\), and both action relations
\[
 G_{0i}=\sum_j\kappa_jG_{ji},\qquad
 G_{00}=\sum_{i,j}\kappa_i\kappa_jG_{ij}.
\]
An exact computational realization uses
\begin{equation}
 M(k)=\frac{\sum_j(k_j-\delta_{ij})G_{ij}
                     M(k-e_i-e_j)}{|k|+2},
 \quad |k|>0\ {\rm even},\ k_i>0,
 \qquad\text{\workstar}
 \label{eq:middle-gram}
\end{equation}
with \(M(0)=1\) and odd moments zero. A term of zero multiplicity is
omitted before its reduced multi-index is formed, so no negative
exponent is evaluated. Gaussian integration by parts
followed by radial separation derives this recurrence; its project
contribution is the complete action/observable implementation and audit,
not a new sphere-moment principle.

Integrating one actual surrounding link changes the problem again.
For \(x=\Tr U/2\), \(y=\Tr V/2\), and
\(w=\Tr(UV^\dagger)/2\), character convolution gives
\[
 \mathbb E[x^ay^bw^c]
 =2^{-a-b-c}\sum_n\frac{m(a,n)m(b,n)m(c,n)}{n+1},
 \quad
 m(a,n)=\binom a{(a-n)/2}-\binom a{(a-n)/2-1}.
\]
In particular \(\mathbb E[xyw]=1/16\) and
\(\mathbb E[x^2y^2w^2]=1/48\).
The actual six affected faces give \(S_2=3x+2y+w\), with
\(O=(4x^2-1)^3(4w^2-1)/81\). Its normalized expectation at
\(\kappa=1/64\) is enclosed near \(1.95552781\,10^{-6}\).
The surrounding marginal is the semicircle measure multiplied by
\(e^{2\kappa y}Z_U(y)\); averaging normalized inner expectations against
bare Haar discards both weights.

Unfreezing a third actual link gives
\(S_3=3x+y+z+w+t\), where \(t=\Tr(VW^\dagger)/2\).
The same \(V\) occurs in both branches:
\(\mathbb E[xzwt]=1/64\), whereas resampling it independently gives zero.
Round19 integrated all required monomials through degree16, retained the
full Taylor tails, and obtained a continuous result for
\(F(\kappa)=\mathbb E[Oe^{\kappa S_3}]/\mathbb E[e^{\kappa S_3}]\):
\begin{equation}
 F(0)=0,\qquad F(\kappa)\ge\kappa^2/2048
 \quad(|\kappa|\le1/8),\qquad
 F(\kappa)>F(-\kappa)\quad(0<\kappa\le1/64).
 \quad\text{\workstar}
 \label{eq:middle-static-sign}
\end{equation}
Here \(N_2=1/324\), \(N_3=13/1296\) for the numerator.
Subtracting absolute higher-order coefficients and its complete tail
bounds \(N/\kappa^2\); an upper partition bound gives the displayed
direction of the quotient. The sign comparison uses
\(N(\kappa)Z(-\kappa)-N(-\kappa)Z(\kappa)\), whose leading term is
\(13\kappa^3/648\). This conditional three-link result still fixes thirty
links. Full surrounding integration and physical dynamical matching remain
separate tasks. Sources:
\repo{research/round17/forward/c1/report.md},
\repo{research/round18/forward/c1/report.md},
\repo{research/round18/forward/c2/report.md},
\repo{research/round19/forward/c1/report.md},
\repo{research/round19/forward/c2/report.md}.

\section{Physical spectral bounds and an explicit infinite-volume exception}
\label{sec:middle-spectral}

\subsection{Finite untruncated graph bounds and complete complements \workstar}
\label{sec:middle-finite-gap}

On a connected finite simple graph with Gauss law at every vertex and no
external charges, a nonconstant spin network has no degree-one active
vertex. Its support therefore contains a cycle. Every occupied link costs
at least \(3\alpha/4\); a fundamental character on a shortest cycle
attains the resulting energy. For graph girth \(g\),
\(\delta_G=3\alpha g/4\) is the exact free physical gap. A forest has only
the vacuum physical sector, not a numerical first excitation.

Let \(L=\sum_p|\lambda_p|\). Bounded multiplication gives
\(E_1(H)\ge\delta_G-L\), while the vacuum trial gives \(E_0(H)\le0\).
For square graphs, the orthonormal trial vectors \(1,\chi_p\) have electric
energies \(0,3\alpha\) and magnetic vacuum-to-face entries
\(-\lambda_p/2\). The relevant third Haar moments vanish. Diagonalizing
only this arrowhead yields a ground-energy \emph{upper} bound. Combining
it with the independent full-space excited-energy bound gives
\begin{equation}
 \Delta(H)\ge
 \frac{\delta+\sqrt{\delta^2+\sum_p\lambda_p^2}}2-L,
 \qquad \delta=3\alpha.
 \qquad\text{\workstar}
 \label{eq:middle-arrow}
\end{equation}
This is not a Ritz-gap lower bound. On the six-face cube, it expands the
common-magnitude sufficient range from
\(r=|\lambda|/\alpha<1/2\) to \(r<12/23\), and gives
\((\sqrt{21/2}-3)\alpha/2>0\) at \(r=1/2\).
On the eleven-face two-cube graph it gives \(r<12/43\).
At that latter boundary an actual shared-face adjoint state
\(\chi_2(U_s)=4x_s^2-1\), of electric energy \(8\alpha\), repairs the
zero margin. With
\(v_0=1+\sum_p\chi_p/22\), \(v_\eta=v_0+\eta\chi_2(U_s)\),
\[
 \frac{\Delta}{\alpha}\ge
 \frac{(6/473)\eta-(347/43)\eta^2}{45/44+\eta^2}>0
 \quad(0<\eta<6/3817).
 \qquad\text{\workstar}
\]
At \(\eta=3/3817\) the exact bound is \(1388/284767457\).
This amplitude maximizes the quadratic numerator, not the quotient.

The next improvement required a complete complement, rather than a larger
trial space alone. Round18 classified every physical state below
\(9\alpha/2\): only the vacuum and eleven fundamental faces occur.
For their projector \(P\), \(QH_0Q\ge(9\alpha/2)Q\), attained by a
six-link rectangle. The full cross Gram is
\[
 \begin{gathered}
 (QVP)^*(QVP)=PV^2P-(PVP)^2,\qquad C_{00}=C_{0p}=0,\\
 C_{pp}=\tfrac14\sum_q\lambda_q^2,\qquad
 C_{pq}=\lambda_p\lambda_q/4\quad(p\ne q).
 \end{gathered}
\]
Both terms in the subtraction matter: \(V\Omega\) already lies in \(P\).
For \(|\lambda_p|\le r\alpha\),
\(\norm{QVP}^2\le21\alpha^2r^2/4\).
On the codimension-one space \(\Omega^\perp\), comparison with
\[
 \alpha\begin{pmatrix}3&-\sqrt{21}\,r/2\\
 -\sqrt{21}\,r/2&9/2-11r\end{pmatrix}
\]
gives the full \(E_1\) lower bound. At \(r=3/8\), it proves
\(\Delta/\alpha\ge(27-3\sqrt{70})/16>0\) for the entire signed box.

Round19 raised the strict electric projector to \(E_{\rm el}<6\alpha\).
Its complete dimension is 48: vacuum, eleven four-cycles and thirty-six
six-cycles. At the excluded threshold \(6\alpha\), 99 support/label
assignments represent 107 physical channels, because eight assignments
have intertwiner multiplicity two. Counting supports alone would miss
this distinction. The full 48-channel cross Gram gives
\(\norm{QVR}^2\le19\alpha^2r^2/4\), with \(R=P-P_\Omega\), while
\(RHR\ge\alpha(3-2r)\) and \(QHQ\ge\alpha(6-11r)\). Hence
\begin{equation}
 \frac{\Delta(H)}{\alpha}\ge
 R(r)=\frac{9-13r-\sqrt{100r^2-54r+9}}2>0,
 \quad
 0\le r<\frac{30-2\sqrt{87}}{23}.
 \qquad\text{\workstar}
 \label{eq:middle-48gap}
\end{equation}
In particular \(R(3/8)=(33-6\sqrt5)/16\) and \(R(7/16)=19/32\).
The failed \(r=1/2\) estimate is a failure of this row-envelope
certificate, not a proved gap closure.

These calculations yield full untruncated finite-graph theorems because
the omitted physical spectrum is bounded analytically. They do not give a
dense volume-uniform threshold. Under the usual lattice normalization
\(\alpha=g_H^2/(2a)\), \(\lambda=2/(g_H^2a)\), one has
\(r=4/g_H^4\), so fixed small-\(r\) regimes do not follow a
weak-bare-coupling path \(g_H\to0\) \cite{bauer2023basis}.
Sources:
\repo{research/round15/forward/C2/derivation.md},
\repo{research/round17/forward/b2/report.md},
\repo{research/round18/forward/b1/report.md},
\repo{research/round19/forward/b1/report.md},
\repo{research/round19/forward/b2/report.md}.

\subsection{Dressed strips, literal boundaries, and summable perturbations \workstar}
\label{sec:middle-strips}

A different route obtains volume-uniform results by an explicit restriction
of the interaction family. Pairwise link-disjoint nonzero plaquettes
factorize the \emph{full} link Hilbert space into four-link blocks and free
links. The four-link free gap before Gauss restriction is \(3\alpha/4\),
not the physical loop energy \(3\alpha\). With
\(\max|\lambda_p|/\alpha\le\rho<3/4\), min--max and zero vacuum mean give
\(\Delta\ge\alpha(3/4-\rho)\). Each unique block ground is gauge invariant:
its gauge action is a continuous one-dimensional representation of a finite
product of \(SU(2)\), hence trivial. The physical restriction contains the
same ground and inherits the lower bound without factorizing the physical
Hilbert space.

Three consecutive square faces form a ten-link strip. Dress its two
link-disjoint ends first and add the middle bridge with coefficient \(\mu\).
The bridge contains a free Haar link in this reference, so its mean is
zero and its squared norm on the reference ground is \(1/4\).
Thus
\[
 \Delta_{\rm strip}\ge
 \alpha\!\left(\tfrac34-
 \max(|\lambda_L|,|\lambda_R|)/\alpha\right)-|\mu|
 \ge\alpha/8
 \quad
 \begin{cases}
 |\lambda_L|,|\lambda_R|\le\alpha/2,\\
 |\mu|\le\alpha/8 .
 \end{cases}
 \qquad\text{\workstar}
\]
The zero-mean premise saves one perturbation norm; a generic bounded
perturbation would lose two.

Complete strips anchored at \((4i,2j,z)\) are link-disjoint. Round18's
finite-box rule initially classified incomplete boundary strips as
remainder faces; a fixed face could consequently change coefficient when
the box grew. Round19 replaced that rule by literal restrictions of a
single infinite assignment. All possible clipped components
\(L,M,R,LM,MR,LMR\) retain the same bound \(\alpha/8\), using the bare
vacuum for proper clips and dressed ends for the complete strip.

In the incomplete tensor product of complete-strip grounds and free Haar
links, define
\[
 H_{\rm ref}=\sum_C(H_C-E_C)+\sum_{e\ {\rm free}}\alpha C_e,
 \qquad H_{\rm ref}\ge\delta(1-P_\Omega),\quad\delta=\alpha/8.
\]
This is the closed nonnegative form sum on the finite-excitation core.
Give every omitted face of anchor \((x,y,z)\) the coefficient
\(\nu_f=\alpha\tau\,2^{-x-y-z}/24\).
The exact geometric sums are
\begin{equation}
 \sum_{\rm all}w_f=1,\quad
 \sum_{\rm selected}w_f=28/135,\quad
 \sum_{\rm omitted}w_f=107/135,
 \qquad
 \norm V\le\beta=\alpha|\tau|\,107/135.
 \qquad\text{\workstar}
 \label{eq:middle-dyadic}
\end{equation}
Every omitted face has a free Haar link, so
\(\langle\Omega,V\Omega\rangle=0\).
On \(\Omega^\perp\), \(H_{\rm ref}+V\ge\delta-\beta\), whereas the
vacuum trial energy is zero. If \(\beta<\delta\), a spectral subspace
below \(\delta-\beta\) is nonempty and at most one-dimensional. This
proves an isolated simple ground without invoking infinite-volume compact
resolvent. At \(|\tau|=1/64\),
\begin{equation}
 \Delta\ge973\alpha/8640.
 \qquad\text{\workstar}
 \label{eq:middle-summable-gap}
\end{equation}
All actual face coefficients can be nonzero, but the omitted ones decay
with position. This is a fully supported inhomogeneous exception.

Round20 completed each finite retained perturbation to whole reference
factors. If \(t_L\downarrow0\) is its exact omitted dyadic tail and
\(\epsilon_L=\alpha|\tau|t_L\), the correct lift is
\(H_L=K_L\otimes I+I\otimes H_{{\rm ref,out}}\).
Standard bounded-perturbation identities \cite{teschl2014methods} give
\[
 \norm{(H-z)^{-1}-(H_L-z)^{-1}}
 \le\epsilon_L/|\operatorname{Im}z|^2,\quad
 |e-e_L|\le\epsilon_L,\quad
 \norm{P-P_L}\le\min(1,\epsilon_L/g_L),
\]
where \(g_L=\delta-\norm{V_L}_{\rm budget}>0\).
The projector bound follows by applying the complementary inverse of
\(H_L-e\) to its residual on the true ground. A zero exterior generator
would leave escaping free excitations at zero and fails norm-resolvent
convergence even at zero perturbation.

Aligned literal vertex boxes \(N=4m+3\) coincide with full factors after
subtracting the explicit reference ground energy
\(c_N=\sum_{C\subset\Lambda_N}E_C\). Other upper boundary phases admit
local-state convergence: complete the support of both a fixed observable
and an interior perturbation, compare their identical central ground
states, then bound each omitted tail. This gives
\[
 |\omega_{\rm box}(A)-\omega_\infty(A)|
 \le4\norm A\min(1,\epsilon_M/g_M)
\]
once the box contains that common completion. It does not establish
arbitrary-phase global vector or norm-resolvent convergence.
Sources: \repo{research/round18/forward/a1/report.md},
\repo{research/round19/forward/a1/report.md},
\repo{research/round19/forward/a2/report.md},
\repo{research/round20/forward/d1/report.md},
\repo{research/round20/forward/e2/report.md}.

\subsection{Physical GNS space and a nonzero Wilson fluctuation \workstar}
\label{sec:middle-gns}

For the dyadic model the full bounded local-factor algebra is irreducible
in the chosen incomplete product. Finite-factor vacuum projections
converge strongly to \(P_\Omega\); an operator in the commutant is scalar
on \(\Omega\), then on its dense local rank-one orbit. Every nonzero
ground vector \(\Psi\) is consequently cyclic for that full algebra.
The GNS null quotient maps isometrically by \([A]\mapsto A\Psi\).
Its energy generator is \(H-e\), since this subtraction fixes the cyclic
vacuum, even though scalar shifts cancel in observable conjugation.

The physical algebra requires an additional argument. Let
\(\mathcal H_{\rm inv}\) be the vectors invariant under every finite-support
vertex gauge transformation, and \(\mathcal A_{\rm phys}\) the norm closure
of bounded invariant local operators. For any invariant \(\xi\), approximate
it by \(A_n\Psi\) using full-algebra cyclicity. Haar-average \(A_n\) over
its finitely many endpoint groups. The weak operator average preserves
local support and, on \(\Psi\), is a strong vector integral satisfying
\(\norm{\mathcal E(A_n)\Psi-\xi}\le\norm{A_n\Psi-\xi}\).
Therefore
\begin{equation}
 \overline{\mathcal A_{\rm phys}\Psi}=\mathcal H_{\rm inv},
 \qquad
 K_{\rm phys}=(H-e)|_{\mathcal H_{\rm inv}},
 \qquad
 \operatorname{spec}K_{\rm phys}\subset
 \{0\}\cup[973\alpha/8640,\infty).
 \qquad\text{\workstar}
 \label{eq:middle-gns}
\end{equation}
The restriction is self-adjoint on \(D(H)\cap\mathcal H_{\rm inv}\);
gauge-commuting spectral cutoffs prove density of that domain.

This sector contains a nonvacuum vector. For the actual omitted origin
\(xz\) Wilson multiplier \(W\), the reference mean is zero and variance
\(1/4\). Two free Haar links per omitted face make distinct-face
covariances vanish. The resulting residual and pure-state trace distance
give, throughout \(|\tau|\le1/64\),
\[
 \begin{gathered}
 d^2=\norm{P_\Psi-P_\Omega}^2
 \le2817/64377572<1/22500,\\
 \operatorname{Var}_\Psi W\ge\tfrac14-2d-4d^2
 \ge5321/22500>\tfrac15 .
 \end{gathered}
 \qquad\text{\workstar}
\]
For \(\xi=(W-\omega(W))\Psi\), spectral calculus now gives the nonzero
imaginary-time correlation
\[
 0<\langle\xi,e^{-tK_{\rm phys}/\hbar}\xi\rangle
 \le\operatorname{Var}_\Psi(W)e^{-973\alpha t/(8640\hbar)}
 \quad(t\ge0).
\]
A lower spectral bound supplies this upper decay bound, not a lower
exponential estimate. It also supplies no real-time magnitude decay.
The result concerns the specified summable representation; equality to a
homogeneous or continuum representation has not been proved.
Sources: \repo{research/round20/forward/g2/report.md},
\repo{research/round21/forward/j1/report.md},
\repo{research/round21/forward/j2/report.md}.

\section{Locality, profile limits, and stationary observations}
\label{sec:middle-dynamics}

The Round13 locality study applied the Lieb--Robinson interaction-picture
method to full \(SU(2)\) links with unbounded onsite Casimirs. Common onsite
unitaries preserve support; only bounded plaquette interactions enter
the chain expansion. If \(q_{\rm inc}=2(d_s-1)\),
\(J=\hbar^{-1}\int j(t)\,dt\) is the dimensionless integrated common
magnetic envelope, and \(R\) is the shortest overlapping-plaquette chain
from the observable support \(X\) to a removed boundary plaquette, then
\begin{equation}
 \norm{\tau_{\Lambda'}(A)-\tau_\Lambda(A)}
 \le\norm A\min\left\{2,\frac{|X|}{4}
                   \sum_{k\ge R}\frac{(8q_{\rm inc}J)^k}{k!}\right\}.
 \qquad\text{\workstar}
 \label{eq:middle-lr}
\end{equation}
There are at most \(|X|q_{\rm inc}(4q_{\rm inc})^{k-1}\) chains;
each commutator costs two and ordered times give \(J^k/k!\).
Unitary resummation is essential: unrestricted growing-support word
counting would not give this chain bound. The factorial tail tends to zero
at fixed \(J\) as the boundary recedes. This proves a fixed-spacing
volume comparison, with no ground-state comparison or continuum velocity.
The underlying strong-operator framework permits infinite-dimensional
sites \cite{nachtergaele2014dynamics}.

The summable model admits a separate direct Duhamel construction.
Completing the local observable and interior perturbation to common
factors gives all-phase finite-box dynamics convergence at fixed time,
with error \(4\norm A\,|t|\epsilon_M/\hbar\).
Strongly continuous implementing unitaries do not imply point-norm
time continuity for every bounded local operator. On a free link the
character shift \(\chi_n\mapsto\chi_{n+1}\) acquires phases with arbitrarily
large frequency; at times tending to zero its norm difference is exactly
two. This limitation was retained when identifying the GNS generator.

\subsection{An exact spatial budget and its endpoint \workstar}
\label{sec:middle-profile}

Replace only the omitted weights by \(w_f(q)=q^{x+y+z}/24\), \(0<q<1\).
Summing all orientations and subtracting selected strips gives
\begin{equation}
 \mathcal B(q)=
 \frac{2+5q+5q^2+6q^3+3q^4}
 {24(1-q)^3(1+q)^2(1+q^2)},\quad
 \mathcal B(1/2)=107/135,\quad
 (1-q)^3\mathcal B(q)\longrightarrow7/64.
 \qquad\text{\workstar}
 \label{eq:middle-profile}
\end{equation}
The same bounded perturbation proof gives
\(\Delta\ge\alpha[1/8-|\tau|\mathcal B(q)]\).
At \(\tau=1/64\) its sufficient boundary lies between \(0.764003\)
and \(0.764004\). This is a boundary of the certificate, not a measured
critical parameter.

Fix \(0<\eta<1\) and follow the canonical constant-budget ray
\(\tau_q=\eta/[8\mathcal B(q)]\).
Each omitted face has two free reference links, and distinct squares
share at most one link. Conditional Haar integration therefore gives
\(\langle x_fx_g\rangle_\Omega=\delta_{fg}/4\), despite possible
entanglement within strip factors. Hence
\begin{equation}
 \sigma_q^2:=\norm{V_q\Omega}^2
 =\frac{\alpha^2\tau_q^2}{96}\mathcal B(q^2),\quad
 \norm{V_q}=\alpha\eta/8,\quad
 \frac{\sigma_q^2}{\alpha^2(1-q)^3}\longrightarrow\frac{\eta^2}{5376}.
 \qquad\text{\workstar}
 \label{eq:middle-variance}
\end{equation}
The norm equality follows by localizing finitely many links near identity
and then controlling the remaining absolute tail; it does not use a
delta function on an infinite identity configuration.
With \(\bar g=\alpha(1-\eta)/8\), the excited threshold gives
\[
 \norm{P_q-P_\Omega}\le\sigma_q/\bar g,\qquad
 -\sigma_q^2/\bar g\le e_q\le0.
\]
Thus the global ground projection returns to the selected-strip reference
although the perturbation norm stays nonzero. Every fixed omitted
coefficient tends to zero. This route cannot retain a nonzero homogeneous
omitted coupling at fixed \(\alpha/E_\star\).

For a bounded operator \(A\) on finitely many complete factors \(F\),
only finitely many omitted faces meet all their links \(E(F)\).
Writing their budget as \(D_F(q)\le|E(F)|/6\),
\[
 \norm{V_qA\Omega}
 \le\norm A[\sigma_q+2\alpha\tau_qD_F(q)]\longrightarrow0.
\]
The local rank-one orbit is dense and \(\norm{V_q}\) is uniformly bounded,
so \(V_q\to0\) strongly. The resolvent identity and strong Duhamel formula
give strong-resolvent convergence and strong unitary convergence uniformly
on compact time intervals. Neither norm-resolvent convergence nor its
failure follows merely from the nonvanishing perturbation norm.

\subsection{Stationarity improves the time window \workstar}
\label{sec:middle-windows}

Round21 first compared connected correlations using a vector residual,
obtaining sufficient windows
\(T_q=(\hbar/\alpha)C(1-q)^{-\gamma}\), \(\gamma<3/2\).
Round22 exploited stationarity before estimating. For bounded complex
local observables,
\[
 C_q^{A,B}(t)=
 \omega_q(A^*\beta_q(t)(B))-\omega_q(A^*)\omega_q(B),
 \qquad \beta_q(t)(B)=e^{-itH_q/\hbar}Be^{itH_q/\hbar}.
\]
Both scalar energy phases cancel in this Heisenberg expression.
The bounded cocycle proves the strong integral identity
\[
 \beta_q(t)(B)-\beta_0(t)(B)
 =-\frac{i}{\hbar}\int_0^t
 \beta_q(t-s)([V_q,\beta_0(s)(B)])\,ds.
\]
No differentiation of an arbitrary bounded \(B\) on an unbounded domain
is needed. Reference evolution stays on the same complete factors,
so \(\norm{[V_q,\beta_0(s)(B)]}\le
2\alpha\tau_q\norm B\,D_{F_B}(q)\) for every \(s\).
Changing the state costs \(2d_q\) in the first expectation and \(4d_q\)
in the two mean factors. Consequently
\begin{equation}
 |C_q^{A,B}(t)-C_0^{A,B}(t)|
 \le\norm A\norm B
 \left[6d_q+\frac{2\alpha\tau_q|t|}{\hbar}D_{F_B}(q)\right].
 \qquad\text{\workstar}
 \label{eq:middle-stationary}
\end{equation}
Since \(d_q=O((1-q)^{3/2})\) and \(\tau_q=O((1-q)^3)\), this includes
every \(\gamma<3\), with sufficient rate
\(O((1-q)^{\min(3/2,\,3-\gamma)})\).
The old \(\gamma=3/2\) endpoint is covered. Failure of the new upper
certificate at \(\gamma=3\) does not establish actual nonconvergence.

For moving origin-anchored tail-block cubes of side \(L\), all outgoing
links are retained: there are \(24L^3\) links and \(21L^3\) incident
omitted faces. Downward closure in the positive orthant makes incidence
equivalent to having the face anchor inside that tail rectangle. Thus
\begin{equation}
 D_L(q)=\mathcal B(q)(1-q^{4L})(1-q^{2L})(1-q^L).
 \qquad\text{\workstar}
 \label{eq:middle-moving}
\end{equation}
For \(L=\max(1,\lfloor\ell(1-q)^{-\beta}\rfloor)\),
\(\ell>0\), \(\beta\ge0\), and time windows
\(|t|\le (\hbar/\alpha)C(1-q)^{-\gamma}\) with \(C>0\), \(\gamma\ge0\),
the nonnegative certificate in \eqref{eq:middle-stationary} vanishes
exactly when
\(\beta<1\) and \(\gamma<3(1-\beta)\).
This necessity concerns that certificate only. At \(\beta=0\), its
constant uses \(L_0=\max(1,\lfloor\ell\rfloor)\), not \(\ell\).
Logarithmically shortened fixed-support \(\gamma=3\) windows also pass.
Translated rectangles need new incoming-face counts. Actual long-time
endpoint behavior remains open.
Sources: \repo{research/round13/locality/locality.md},
\repo{research/round20/forward/h2/report.md},
\repo{research/round21/forward/m1/report.md},
\repo{research/round21/forward/m2/report.md},
\repo{research/round22/reverse/n1/report.md},
\repo{research/round22/reverse/n2/report.md}.

\section{Static information, clocks, and exact physical reduction}
\label{sec:middle-inference}

\subsection{Conditional dynamics and finite-rate identifiability \workstar}
\label{sec:middle-clock}

On the conditional three-link space with
\(\rho_\kappa\propto e^{\kappa S_3}\), a separately specified reversible
generator is
\[
 A_{\kappa,m}=-\rho_\kappa^{-1}
       \operatorname{div}(\rho_\kappa m\nabla),\qquad H=cA_{\kappa,m},
 \quad c>0.
\]
For smooth uniformly positive mobility \(m\), its closed form has domain
\(H^1\), operator domain \(H^2\), and constant ground. This standard
divergence-form construction preserves the same static law for every \(c\)
and \(m\), demonstrating why static moments cannot fix a clock.
The actual shared-middle gradient produces
\[
 \Gamma(S_3)=\tfrac14\{15+8y+2x+2z+2r
 -(3x+w)^2-(y+w+t)^2-(z+t)^2\},\quad
 r=\Tr(UW^\dagger)/2.
\]
The mixed term involving \(r-wt\) disappears under an incorrect resampling
of \(V\). Independent optimization of \(U,W\) at fixed \(y\) proves
\(\min S_3=-5\), \(\max S_3=7\).
Bounded-density comparison with the Haar Casimir gap then gives
\(\operatorname{gap}(cA_{\kappa,1})\ge(3c/4)e^{-12|\kappa|}
\ge c/6\) for \(|\kappa|\le1/8\).
This is a theorem for the declared conditional generator, not a derived
lattice Hamiltonian.

For \(m=1+\zeta x\), \(|\zeta|<1\), the measurable initial rate would be
\(r_f=-\hbar C_f'(0+)=c\int m\Gamma(f)\rho_\kappa\).
At known \(\kappa=0\), the pair \(x,x+x^2\) gives
\[
 r_x=3c/16,\quad r_{x+x^2}=c(5/16+\zeta/8),\quad
 c=16r_x/3,\quad\zeta=8r_{x+x^2}/c-5/2.
 \qquad\text{\workstar}
\]
The pair \(x,x^2\) is parity-blind to \(\zeta\).
For known \(|\kappa|\le1/8\), choosing \(g=x+x^2/4\) proves
a uniform determinant lower bound \(D\ge1/7168\). Indeed, if
\(G(x)=(1-x^2)/4\) and \(X,X'\) are independent copies of the actual
one-coordinate marginal, symmetrization gives, for \(g=x+ax^2\),
\[
 D=2a\,\mathbb E\{G(X)G(X')(X-X')^2[1+a(X+X')]\}>0
 \quad(0<a<1/2).
\]
Comparison to Haar yields
\(D\ge3a(1-2a)e^{-24|\kappa|}/128\); \(a=1/4\) and \(e^3<21\)
give the stated rational floor.

These inverses identify only their declared finite family. At Haar,
\(p_3=x^3-3x/8=C_3^{(2)}(x)/32\) is invisible to both rates.
Adding \(h=x^2+x^3\) identifies \(m=1+\zeta x+\xi p_3\) through
\begin{equation}
 \begin{pmatrix}r_x\\r_{x+x^2/4}\\r_{x^2+x^3}\end{pmatrix}
 =
 \begin{pmatrix}
 3/16&0&0\\25/128&1/32&0\\59/256&9/64&9/512
 \end{pmatrix}
 \begin{pmatrix}c\\c\zeta\\c\xi\end{pmatrix},
 \qquad \det M=27/262144.
 \qquad\text{\workstar}
 \label{eq:middle-rates}
\end{equation}
The box \(|\zeta|,|\xi|\le1/4\) guarantees \(m\ge13/32\), hence
gap at least \(13c/192\) on the stated \(\kappa\) range.
The Gegenbauer polynomial itself is classical
\cite{nistDLMFGegenbauer}; its blind and detecting roles here are the
project calculation.

More generally \(N\) initial rates impose \(N\) linear conditions on
a polynomial with \(N+1\) coefficients. A nonzero null polynomial,
normalized by its coefficient absolute sum, has sup norm at most one.
It gives a whole positive interval of distinct mobilities with identical
rates at fixed density and \(c\). The explicit
\(p_5=x^5-5x^3/6+x/8=C_5^{(2)}(x)/192\) preserves all three rates above,
but changes the \(x\)-correlation curvature to
\[
 C_x''(0+)=\frac{c^2}{\hbar^2}
          \left(\frac9{64}+\frac{\epsilon^2}{1024}\right)
 \quad\text{for }m=1+\epsilon p_5.
 \qquad\text{\workstar}
\]
Thus the obstruction concerns finitely many initial slopes, not complete
time curves. No physical rates were measured; the inversions were tested
on synthetic exact inputs. Sources:
\repo{research/round20/forward/f2/report.md},
\repo{research/round21/forward/k2/report.md},
\repo{research/round21/forward/l1/report.md},
\repo{research/round21/forward/l2/report.md}.

\subsection{A failed section and a valid maximal-tree map \workstar}
\label{sec:middle-tree}

Round22 tested physical matching on the actual 18-vertex, 33-link graph
at zero magnetic coupling. Fixing thirty links to identity leaves the
old \(U,V,W\) coordinates, but fixes thirteen genuine cycles in addition
to tree gauge variables. It is not a Hilbert-space reduction.
If \(Q_0\) is a face fixed to identity, then
\(f_n=\chi_n(Q_0)/(n+1)\) satisfies
\(\norm{f_n}_2=1/(n+1)\to0\), whereas its section equals one.
The restriction is nonclosable. Three edge-disjoint such faces give an
analogous sequence converging even in the electric graph norm.
Moreover \(H_E x_0=3\alpha x_0\) before restriction, but the conditional
generator annihilates its constant image. A scalar clock cannot repair
this core defect.

A designated slope fit for a physical square that restricts to \(x\)
gives \(c=4\alpha\). A held-out square restricting to
\(w=\Tr(UV^\dagger)/2\) has physical correlation
\(e^{-3\alpha t/\hbar}/4\), while the fitted conditional prediction is
\(e^{-6\alpha t/\hbar}/4\). The disagreement is strict at every
finite positive time.

The repair uses a genuine maximal tree
\(T=\{0,\ldots,15,26\}\), root \((0,2,0)\), and all sixteen chord
holonomies \(L_e=h_{s(e)}g_eh_{t(e)}^{-1}\).
At fixed tree variables, left/right Haar invariance gives product Haar
on the chords. Gauge-invariant functions become simultaneous-conjugation
invariants on these sixteen variables, yielding a unitary \(J_{16}\).
The three-chord isometry \(J_3\) embeds functions of \(L_{28},L_{27},L_{24}\);
its adjoint integrates the other thirteen chords, rather than evaluating
them at identity. These are the standard maximal-tree principles
\cite{burbano2024gauge}, here realized on the complete signed graph.

Every tree derivative must also be transported. If cutting tree edge \(q\)
leaves subtree \(S_q\), the induced vector field is
\[
 D_{q,a}=\epsilon_q\left(
 \sum_{s(e)\in S_q}\ell_{e,a}
 -\sum_{t(e)\in S_q}\rho_{e,a}\right).
\]
Both endpoint occurrences remain when a chord lies inside \(S_q\).
The full generator is
\(-\alpha\sum_a(\sum_e\ell_{e,a}^2+\sum_qD_{q,a}^2)\).
Its quadratic form is a sum of squared gradients, bounded above and below
by positive multiples of the product Casimir form. On finite
Peter--Weyl blocks these comparisons give equivalence of \(H^1\) forms
and \(H^2\) operator domains and identify the closed operators.
On the selected three chords the exact reduction is
\begin{equation}
 H_{\rm eff}/\alpha=
 3C_U+3C_V+C_W+C^\ell_{VW}+C^\rho_{VW}
                  +C^\ell_{UVW}+2C^\rho_{UVW}.
 \qquad\text{\workstar}
 \label{eq:middle-effective}
\end{equation}
Here \(C^\ell_A=-\sum_a(\sum_{j\in A}\ell_{j,a})^2\), and similarly
on the right. The image reduces \(H_E\) because its orthogonal projector
is Haar integration over unused original chord links and commutes with
each link heat semigroup. This proves exact electric semigroup
intertwining, not merely a formal compression.

The new \(x,y,z\) observables complete to physical loops of lengths
six, eight, six, respectively, so their electric energies are
\((9/2)\alpha,6\alpha,(9/2)\alpha\).
They are different completions from the old section's training square.
The prescribed \(x\)-slope now gives \(c=6\alpha\); its reserved
curvature forces \(\zeta=0\). The held-out \(y\)-curve still differs:
the physical exponent is \(6\alpha\), the conditional one \(9\alpha/2\).
Thus exact Haar reduction succeeds while matching to the scalar-clock
affine-mobility family fails. Sources:
\repo{research/round22/reverse/p1/report.md},
\repo{research/round22/reverse/p2/report.md}.

\subsection{Magnetic leakage, projected memory, and exact channels \workstar}
\label{sec:middle-memory}

Keep this tree and add all twenty magnetic faces:
\(H_\lambda=H_E+\alpha\lambda V\), \(V=\sum_{p=0}^{19}(1-x_p)\),
\(\lambda\ge0\). This \(\lambda\) is dimensionless, unlike the physical
\(\lambda_p\) used earlier. With \(P=J_3J_3^*\), only one face survives
conditional averaging, so \(J_3^*VJ_3=20-\Tr(VW^\dagger)/2\).
Three other faces share the same omitted chord. Their nonzero conditional
cross moments give
\begin{equation}
 \begin{gathered}
 B_\lambda^*B_\lambda=(\alpha\lambda)^2M_d,\qquad
 B_\lambda=QH_\lambda J_3,\\
 d=\tfrac{19}{4}+\tfrac12(r+x+z)
   =4+\tfrac14\norm{u+w+e_0}^2\in[4,25/4].
 \end{gathered}
 \qquad\text{\workstar}
 \label{eq:middle-leakage}
\end{equation}
Here \(r=\Tr(UW^\dagger)/2\), and \(u,w\) are unit quaternions.
Hence every nonzero selected vector leaks at nonzero \(\lambda\);
\(\norm{B_\lambda}=5\alpha\lambda/2\).
Replacing \(d\) by its mean \(19/4\) loses the correlated face terms.
Also \(H_\lambda1\) is nonconstant, with centered squared residual
\(5(\alpha\lambda)^2\); the Haar reference is not an interacting ground.

Set \(A_\lambda=J_3^*H_\lambda J_3\) and
\(T_\lambda(t)=J_3^*e^{-tH_\lambda/\hbar}J_3\).
Two bounded off-diagonal Duhamel insertions prove
\[
 \norm{T_\lambda(t)-e^{-tA_\lambda/\hbar}}
 \le\min\{1,\tfrac{25}{8}(\alpha\lambda t/\hbar)^2\}.
 \qquad\text{\workstar}
\]
The scaled difference tends strongly to \(M_d/2\) as \(t\downarrow0\);
this is not an operator-norm Taylor expansion. It already disproves
autonomy of the selected compression.

The complement has the exact electric lower bound \(3\alpha\), by the
physical girth argument and an omitted square attaining it. Thus
\(C_\lambda=QH_\lambda Q\ge3\alpha\). Eliminating the complementary
component gives the exact memory and resolvent relations
\[
 \hbar u'=-A_\lambda u+\hbar^{-1}
       \int_0^tK_\lambda(t-s)u(s)\,ds,\quad
 K_\lambda(s)=B_\lambda^*e^{-sC_\lambda/\hbar}B_\lambda,
\]
\[
 J_3^*(H_\lambda+z)^{-1}J_3
 =[A_\lambda+z-\Sigma_\lambda(z)]^{-1},\qquad
 \Sigma_\lambda(z)=B_\lambda^*(C_\lambda+z)^{-1}B_\lambda,\quad z>0.
\]
The full return \(u(s)\) remains inside the memory integral.
These block identities are established methods; the actual channel
content and quantified errors are the project contribution.

For the two orthonormal channels \(1,2x\), the leading kernel
\((\alpha\lambda)^2B_1^*e^{-sC_0/\hbar}B_1\) has seven exact electric
energies. The \(11\) weights at energies
\(3,9/2,6,15/2,8,17/2,9\), in units of \(\alpha\), are
\[
 (1/8,\ 1/8,\ 3/8,\ 9/4,\ 9/8,\ 3/8,\ 3/8).
 \qquad\text{\workstar}
\]
The \(00\) weight is \(19/4\) at energy three, and the off-diagonal
weight is \(1/4\) at that same energy. To derive the table, intersect
each face with the actual six-link \(x\) loop. A shared consecutive
path of length \(m>0\) has spin-zero and spin-one components of squared
norms \(1/16,3/16\), with energies
\(3(10-2m)/4\) and that number plus \(2m\). Nine faces have disjoint
edge support; the other counts for \(m=1,2,3\) are six, two, two.
This is an exact finite spectral orbit in the full Hilbert space,
not a representation cutoff.

Since \(\norm V\le40\), bounded Duhamel and the resolvent identity give
full selected-space errors
\begin{equation}
 \begin{split}
 \norm{K_\lambda(s)-K_\lambda^{(2)}(s)}
 &\le250\alpha^3\lambda^3(s/\hbar)e^{-3\alpha s/\hbar},\\
 \norm{\Sigma_\lambda(z)-\Sigma_\lambda^{(2)}(z)}
 &\le\frac{250\alpha^3\lambda^3}{(3\alpha+z)^2}.
 \end{split}
 \qquad\text{\workstar}
 \label{eq:middle-memory-errors}
\end{equation}
These are uniform absolute kernel bounds in delay and uniform
complementary self-energy bounds down to \(z=0\). They do not give uniform
relative late-time accuracy or a bounded full resolvent at zero.
Sources: \repo{research/round22/reverse/q1/report.md},
\repo{research/round22/reverse/q2/report.md}.

\section{Homogeneous stability: explicit progress and retained obstructions}
\label{sec:middle-homogeneous}

\subsection{A matched theorem dictionary and a controlled rotation \workstar}
\label{sec:middle-rotation}

The qualitative homogeneous stability route is independent of the summable
construction. Yarotsky's product-vacuum theorem permits unbounded onsite
operators and infinite-dimensional site spaces, but its range-dependent
constants remain unevaluated \cite{yarotsky2004quasiparticles}.
The Round13 outgoing-link blocking gives normalized perturbation
\(\epsilon\le(4/3)\binom d2\lambda_{\max}/\alpha\).
Boundary padding handles all-contained plaquettes; it does not silently
identify every thermodynamic boundary limit. A sufficiently small,
existential magnetic/electric ratio gives a fixed-spacing gap at least
\(3\alpha/8\). A qubit, two-site algorithmic theorem cannot supply its
missing constants \cite{bravyi2008polynomial}.

The dressed-strip reference admits a sharper dictionary. Each coarse tail
cell contains 24 links: ten strip links and fourteen free links.
There are 24 anchored faces, of which three are selected and 21 omitted.
Their full grouped support is
\(G=\{0,e_x,e_y,e_z\}\). In energy unit \(\delta=\alpha/8\),
\[
 h_b\ge1-P_b,\qquad
 \phi_b=-\frac{\tau}{3}\sum_{f\in O_b}x_f,\qquad
 \norm{\phi_b}=7|\tau|.
 \qquad\text{\workstar}
\]
All 21 omitted faces cross a coarse block boundary. Ignoring crossing
faces would erase the entire homogeneous perturbation.
The matched source theorem gives an existential interval
\[
 |\tau|<\tfrac17\min\{c_1(G),[2c_2(G)]^{-1}\},
 \qquad \Delta>\alpha/16,
\]
including literal finite rectangles through clipped reference sites and
padding. No numerical value of \(\tau\), including \(1/64\), is thereby
certified.

The full magnetic interaction is not purely form-bounded relative to
the undressed vacuum energy. On an actual square,
\(\psi_t=(1+t\chi_p)/\sqrt{1+t^2}\) has electric energy
\(3\alpha t^2/(1+t^2)\) and magnetic expectation
\(-\lambda t/(1+t^2)\); their absolute ratio diverges as \(t\to0\).
The diagonal compression alone does obey the incidence bound
\(|\langle W_{\rm diag}\rangle|\le
(16\lambda_{\max}/3\alpha)\langle H_0\rangle\).
The missing off-diagonal term cannot be discarded.

On one four-site star, conditional Haar orthogonality gives
\(\norm{\phi\Omega}^2=7\tau^2/12\).
Let \(v=\phi\Omega\), \(u=(H_0|_Q)^{-1}v\), and
\(S=|u\rangle\langle\Omega|-|\Omega\rangle\langle u|\).
Then \([S,H_0]=-(|v\rangle\langle\Omega|+\mathrm{adjoint})\).
The finite-rank range lies in \(D(H_0)\); the exponential preserves that
domain. Integrating the bounded commutator gives the exact identity
\[
 e^S(H_0+\phi)e^{-S}=H_0+Q\phi Q+R,\qquad
 \norm R\le(707/60)\tau^2.
 \qquad\text{\workstar}
\]
This instantiates a standard unitary dressing method with the actual
local variance. The remainder is part of the Hamiltonian.

For simultaneous finite-volume rotations, every nonzero ordered
commutator word adds a star meeting the previous union.
\(|G-G|=13\), so there are at most \(13^nn!\) relative words of order
\(n\), each supported on at most \(4+3n\) sites. Rooting at a fixed site
adds a factor at most \(4+3n\).
In the interaction norm
\(\norm\Phi_w=\sup_x\sum_{Y\ni x}2^{|Y|}\norm{\Phi_Y}\),
this yields a convergent all-order remainder majorant.
The admitted common certificate is
\begin{equation}
 \norm R_w\le\frac{25460736}{25}\tau^2,\qquad
 |\tau|\le5/1664,\qquad
 |\langle D\rangle|\le28|\tau|\langle H_0\rangle.
 \qquad\text{\workstar}
 \label{eq:middle-global-rotation}
\end{equation}
The reverse derivation separately obtained the smaller valid constant
\(1970176/5\); it is not substituted into the subsequently frozen common
recurrence. Neither constant is a homogeneous gap threshold.

\subsection{Generated supports, interacting inverses, and the remaining source \workstar}
\label{sec:middle-inverse}

For a residual term on an arbitrary generated support \(Y\), split
\[
 R_Y=c_YI+A_Y+Z_Y,\quad
 A_Y=|Q_YR_Y\Omega_Y\rangle\langle\Omega_Y|+\mathrm{adjoint},
 \quad Z_Y=Q_Y(R_Y-c_YI)Q_Y .
\]
The subtraction inside \(Z_Y\) is necessary when the scalar is \(c_YI\).
The bare homological inverse has norm at most one, and the resulting
diagonal update is bounded by \(2\norm{R_Y}\).
For exponential cardinality weight \(\mu\), an explicit root/intersection
sum gives
\[
 \norm{[\Phi,\Psi]}_{\mu-\ell}
 \le\frac4{e\ell}\norm\Phi_\mu\norm\Psi_\mu,\qquad \ell>0.
\]
Spending \(\ell/n\) per nested commutator and
\(n!\ge(n/e)^n\) bounds the factorial-divided order by
\((4\norm S_\mu/\ell)^n\norm K_\mu\).
If \(r>0\) bounds the residual, \(b\ge0\) bounds the retained diagonal
interaction, and \(\ell>4r\), the reverse majorant is
\begin{equation}
 \begin{gathered}
 r'\le\mathcal F(r,b;\ell):=
       \frac{4r}{\ell-4r}(b+3r/2),\qquad b'\le b+2r,\\
 \mathcal F(r,b;\ell)<r\ \Longleftrightarrow\ \ell>4b+10r .
 \end{gathered}
 \qquad\text{\workstar}
 \label{eq:middle-recurrence}
\end{equation}
The equivalence refers to the displayed scalar majorant, not the actual
residual; \(r=0\) is the separate exact zero-residual case. Any finite
total loss has \(\ell_n\to0\), while this majorant
retains \(b_n\ge448|\tau|>0\). Its positive scalar trajectory eventually
leaves its convergence radius. This proves failure of that certificate,
not of the true rotation algorithm or the physical gap.

An equal-weight commutator counterexample also needs its exact support
convention. In the reverse structured example, regrouping to the actual
support gives the special homological value \(1/64\); retaining indexed
unions instead gives \((9N+16)/(64(3N+1))\).
The post-freeze gate explicitly records this distinction. It is not
legitimate to improve an interaction norm by silently changing its
decomposition.

Round22 then replaced the bare inverse by the actual \emph{initial}
diagonal operator. For finite \(Y\), let
\(G_Y=H_Y+\sum_{b+G\subset Y}Q_b\phi_bQ_b\).
Four-star incidence gives, with \(a=28|\tau|<1\),
\[
 (1-a)H_Y\le G_Y\le(1+a)H_Y,\qquad
 \norm{(G_Y|_{Q_Y})^{-1}}\le(1-a)^{-1}.
\]
The inverse maps into \(D(H_Y)\), so its rank-two skew generator has a
valid domain-preserving exponential. Interior cancellation is exact,
but on the full finite volume
\[
 [S_Y,H_0+D]+A_Y
 =\sum_{\substack{b+G\cap Y\ne\varnothing\\b+G\not\subset Y}}
                          [S_Y,D_b].
\]
For the admitted actual \(SU(2)\) diagnostic source consisting of a free
link character at the origin, the squared boundary residual on the
product vacuum is exactly
\begin{equation}
 \norm{([S_Y,H_0+D]+A_Y)\Omega}^2=7\tau^2/432.
 \qquad\text{\workstar}
 \label{eq:middle-boundary}
\end{equation}
All 21 face contributions and their vanishing distinct-face cross
moments are retained. This source is admissible but is not asserted to
be an actually generated term of the earlier rotation.

Finally, the infinite initial diagonal is constructed as a relative form
sum on the same incomplete product. Let
\(K=H_0^{-1/2}DH_0^{-1/2}\) on \(Q_\Omega\);
\(\norm K\le a<1\). For a vacuum-exterior source \(v\perp\Omega\),
\begin{equation}
 u=H_0^{-1/2}(I+K)^{-1}H_0^{-1/2}v
   =\sum_{n\ge0}H_0^{-1/2}(-K)^nH_0^{-1/2}v .
 \qquad\text{\workstar}
 \label{eq:middle-neumann}
\end{equation}
The form equation proves \(Gu=v\).
Each coefficient lies in the \(n\)-collar \(Y^{(n)}\) with norm at most
\(\norm v\,a^n\); direct finite-support action also gives an
\(H_0\)-graph bound
\(\norm{H_0u_n}\le\norm v\,|Y|(2n-1)^3a^n\) for \(n\ge1\).
Polynomially weighted geometric tails prove Hilbert, half-energy,
\(H_0\)-graph and \(G\)-graph convergence, including explicit
finite-volume collar errors. This response regularity does not prove
global equality \(D(G)=D(H_0)\).

The exact global vacuum rank-two source is \(A_Y\otimes P_{\rm ext}\).
Against the desired local source \(A_Y\otimes I\), the remaining operator
is \(A_Y\otimes Q_{\rm ext}\), of norm one on the specified exterior
character test. Moreover, the prescribed collar interaction certificate
\(\sum_n e^{\mu(n+1)^3}a^n\) diverges for every fixed \(\mu>0\),
\(0<a<1\), although the unweighted response converges.
Thus full local-source cancellation and a useful weighted decomposition
remain missing. These are the precise next requirements for an all-stage
homogeneous stability proof; the later diagonal updates are not covered
by the initial inverse.

The domain and inverse comparisons use established form and
Lie--Schwinger methods \cite{teschl2014methods,delvecchio2021unbounded}.
Their complete source inductions and numerical stability thresholds
were not imported. Sources:
\repo{research/round21/forward/i1/report.md},
\repo{research/round21/forward/i2/report.md},
\repo{research/round22/advisor/o1-gate.json},
\repo{research/round22/advisor/o2-gate.json},
\repo{research/round22/reverse/r1/report.md},
\repo{research/round22/reverse/r2/report.md}.
